% 核心结论：补齐表头至Rx及电压表左端导线，导线在V圆边界终止
% 草图：[真实几何与节点 | 条件、推导、结论]；17cm印宽，最小字号8pt。
% 图契约：伏安法的两种接法与仪表负载效应；图中的端口、几何、箭头和方程对应本题。
% 版面：沿用前三模块的几何图+推导分区；目标印宽17cm，简图至中等密度。
% 布局：左侧真实结构，右侧关键关系；真实电路边不作为架构流程箭头处理。
\documentclass[tikz,border=7pt]{standalone}
\usepackage{fontspec,xeCJK,amsmath,amssymb}
\IfFontExistsTF{Times New Roman}{\setmainfont{Times New Roman}}{\setmainfont{TeX Gyre Termes}}
\IfFontExistsTF{Heiti SC}{\setCJKmainfont[AutoFakeBold=1.5]{Heiti SC}}{\setCJKmainfont[AutoFakeBold=1.5]{FandolHei-Regular}}
\usetikzlibrary{arrows.meta,bending,calc,patterns,positioning,decorations.pathmorphing}
\definecolor{ink}{HTML}{243745}
\definecolor{blue}{HTML}{275DAB}
\definecolor{green}{HTML}{177C63}
\definecolor{red}{HTML}{BD4935}
\definecolor{orange}{HTML}{AA6C16}
\definecolor{purple}{HTML}{715393}
\tikzset{every node/.style={font=\fontsize{9.3}{12.5}\selectfont,text=ink},
arrow/.style={-{Stealth[length=5pt 1.5,width'=0pt 3]},shorten >=2pt,shorten <=1pt,line width=1pt},
arrow short/.style={-{Stealth[length=3pt 1.0,width'=0pt 2.5]},shorten >=1pt,shorten <=0pt,line width=.8pt},
axis/.style={arrow,draw=ink!60,line width=.6pt},
vec/.style={arrow,draw=blue},
wire/.style={draw=ink!75,line width=.8pt},
guide/.style={dashed,draw=ink!45,line width=.5pt},
card/.style={rounded corners=2pt,draw=blue!65,fill=blue!4,align=center,text width=6.5cm,inner xsep=9pt,inner ysep=7pt},
title/.style={font=\fontsize{10.5}{14}\selectfont\bfseries,text width=15cm,align=center},
label/.style={font=\fontsize{8.4}{11}\selectfont,align=center},
note/.style={label,text width=6.1cm}}
\begin{document}
\begin{tikzpicture}[x=1cm,y=1cm]
\path[use as bounding box] (-8,-3.5) rectangle(8,3.65);
\node[title] at (0,3.25) {\quad 伏安法的两种接法与仪表负载效应};
\begin{scope}[shift={(-4.05,0)}]
\draw[wire,fill=white](-1.7,1.35)circle(0.25);\node at(-1.7,1.35){$A$};
\draw[wire,draw=ink](-1.1,1.35)--(0.32499999999999984,1.35) (0.8749999999999999,1.35)--(2.3,1.35);\begin{scope}[shift={(0.5999999999999999,1.35)},rotate=0.0]\draw[draw=ink,fill=white,line width=.8pt](-0.275,-.085)rectangle(0.275,.085);\end{scope}
\node[label]at(0.5999999999999999,1.77){$R_x$};
\draw[wire](-2.6,1.35)--(-1.95,1.35)(-1.45,1.35)--(-1.1,1.35)(2.3,1.35)--(2.6,1.35);\draw[wire,fill=white](0.45,0.40000000000000013)circle(0.25);\node at(0.45,0.40000000000000013){$V$};
\draw[wire](-2.3,1.35)--(-2.3,0.40000000000000013)--(.2,0.40000000000000013)(.7,0.40000000000000013)--(2.3,0.40000000000000013)--(2.3,1.35);\fill[ink] (-2.3,1.35)circle(2pt);
\fill[ink] (2.3,1.35)circle(2pt);
\node[label]at(0,2.1500000000000004){电压表跨接表头与待测电阻};\draw[wire,fill=white](-1.7,-1.1)circle(0.25);\node at(-1.7,-1.1){$A$};
\draw[wire,draw=ink](-1.1,-1.1)--(0.32499999999999984,-1.1) (0.8749999999999999,-1.1)--(2.3,-1.1);\begin{scope}[shift={(0.5999999999999999,-1.1)},rotate=0.0]\draw[draw=ink,fill=white,line width=.8pt](-0.275,-.085)rectangle(0.275,.085);\end{scope}
\node[label]at(0.5999999999999999,-0.6800000000000002){$R_x$};
\draw[wire](-2.6,-1.1)--(-1.95,-1.1)(-1.45,-1.1)--(-1.1,-1.1)(2.3,-1.1)--(2.6,-1.1);\draw[wire,fill=white](0.45,-2.05)circle(0.25);\node at(0.45,-2.05){$V$};
\draw[wire](-1.1,-1.1)--(-1.1,-2.05)--(.2,-2.05)(.7,-2.05)--(2.3,-2.05)--(2.3,-1.1);\fill[ink] (-1.1,-1.1)circle(2pt);
\fill[ink] (2.3,-1.1)circle(2pt);
\node[label]at(0,-0.30000000000000004){电压表只跨接待测电阻};\end{scope}
\node[card,text width=6.5cm,draw=blue!65,fill=blue!4]at(3.9,1.1){上图：$U/I=R_x+R_A$\\电流表内阻使读数偏大};
\node[card,text width=6.5cm,draw=green!65,fill=green!4]at(3.9,-1.1){下图：$U/I=R_x\parallel R_V$\\电压表分流使读数偏小};

\end{tikzpicture}\end{document}